\documentclass[greek]{beamer}
%\usepackage{fontspec}
\usepackage{amsmath,amsthm}
\usepackage{unicode-math}
\usepackage{xltxtra}
\usepackage{graphicx}
\usetheme{CambridgeUS}
\usecolortheme{seagull}
\usepackage{hyperref}
\usepackage{ulem}
\usepackage{xgreek}

\usepackage{pgfpages} 
\usepackage{tikz}
%\setbeameroption{show notes on second screen}
%\setbeameroption{show only notes}

\setsansfont{Calibri}

\usepackage{multicol}

\usepackage{appendixnumberbeamer}

\usepackage{polynom}

\usepackage{pgffor}

\setbeamercovered{transparent}
\beamertemplatenavigationsymbolsempty

\title{Συναρτήσεις}
\subtitle{Όριο Συνάρτησης στο $x_0\in\mathbb{R}$}
\author[Λόλας]{Κωνσταντίνος Λόλας }
\institute[$10^ο$ ΓΕΛ]{$10^ο$ ΓΕΛ Θεσσαλονίκης}
\date{}

\begin{document}

\begin{frame}
      \titlepage
\end{frame}

\section{Θεωρία}
\begin{frame}{Όριο}
      \centering
      \includegraphics[width=0.6\columnwidth]{images/6r7k74}
\end{frame}

\begin{frame}{Όριο}
      \begin{block}{Το αστέρι μας}
            \centering
            $\lim\limits_{x \to x_0}{ f(x) }$ \pause
      \end{block}

      Διαβάζεται ως:
      \begin{itemize}
            \item Το όριο της εφ όταν το χι τείνει στο χιμηδεν \pause
            \item Το όριο της $f$ στο $x_0$ \pause
            \item Όταν το $x$ πάει στο $x_0$, πού πάει η $f$...
      \end{itemize}
\end{frame}

\begin{frame}{Ξεπηδούν οι απορίες}
      \begin{itemize}
            \item Τι σημαίνει πλησιάζω στο $x_0$ \pause
                  \begin{itemize}
                        \item Δημιουργήστε την γραμμή των πραγματικών αριθμών και πλησιάστε στο $x=2$ \pause
                        \item Με πόσους τρόπους μπορείτε να πλησιάσετε
                  \end{itemize}
            \item Τι σημαίνει η $f$ πλησιάζει στο $l$ \pause
            \item Τι σημαίνει οσοδήποτε κοντά
      \end{itemize}
\end{frame}

\begin{frame}{Ας γίνουμε νονοί}
      \begin{block}{Αριστερό πλευρικό όριο}
            \begin{center}
                  $\lim\limits_{x \to x_0^-}{ f(x) }$ \pause
            \end{center}

            Για μια συνάρτηση που ορίζεται σε διάστημα της μορφής $(α,x_0)$ για κατάλληλο $α$
      \end{block}
\end{frame}

\begin{frame}{Ας γίνουμε νονοί}
      \begin{block}{Αριστερό πλευρικό όριο}
            \begin{center}
                  $\lim\limits_{x \to x_0^-}{ f(x) }$ \pause
            \end{center}
      \end{block}
      \begin{block}{Δεξί πλευρικό όριο}
            \begin{center}
                  $\lim\limits_{x \to x_0^+}{ f(x) }$ \pause
            \end{center}

            Για μια συνάρτηση που ορίζεται σε διάστημα της μορφής $(x_0,α)$ για κατάλληλο $α$
      \end{block}
\end{frame}

\begin{frame}{Άρα}
      \begin{block}{Ύπαρξη ορίου}
            \begin{equation*}
                  \lim\limits_{x \to x_0^-}{ f(x) }=λ\iff
                  \begin{cases}
                        \lim\limits_{x \to x_0^-}{ f(x) }=λ\in\mathbb{R} \\
                        \lim\limits_{x \to x_0^+}{ f(x) }=λ\in\mathbb{R} \\
                        \lim\limits_{x \to x_0^-}{ f(x) }=\lim\limits_{x \to x_0^+}{ f(x) }
                  \end{cases}
            \end{equation*}
      \end{block}
\end{frame}

\begin{frame}{Περιπτωσάρα}
      Αν $f(x)=\sqrt{x}$?, ή $f(x)=\ln (-x)$? \pause
      \begin{block}{ }
            Αν μια συνάρτηση ορίζεται μόνο σε διάστημα της μορφής $(α,x_0)$ τότε $\lim\limits_{x \to x_0}{ f(x) }=\lim\limits_{x \to x_0^-}{ f(x) }$
      \end{block} \pause
      Όμοια για $\lim\limits_{x \to x_0}{ f(x) }=\lim\limits_{x \to x_0^+}{ f(x) }$
\end{frame}

\begin{frame}{Ιδιότητες}
      \begin{itemize}
            \item Το όριο στην περίπτωση που υπάρχει είναι μοναδικό \pause
            \item $\lim\limits_{x \to x_0}{ f(x) }=k\iff\lim\limits_{x \to x_0}{ \left( f(x)-k \right) }=0$ \pause
            \item $\lim\limits_{x \to x_0}{ f(x) }=k\iff\lim\limits_{h \to 0}{ f(h+x_0) }=k$
      \end{itemize}
\end{frame}
\begin{frame}{Άρα τι θα κάνουμε?}
      \begin{itemize}
            \item Θα περιγράφουμε \pause
            \item Θα υπολογίζουμε (χωρίς να ξέρουμε γιατί) \pause
            \item Θα χρησιμοποιούμε ιδιότητες και τεχνικές \pause
            \item αλλά και πάλι δεν θα καταλαβαίνουμε
      \end{itemize}
      Ουσιαστικά τα όρια θα τα υπολογίζουμε εντελώς μηχανικά
\end{frame}

\begin{frame}{Επίδειξη}
      Στο διάλλειμα όποιος θέλει μπορεί να μάθει τον υπέρτατο ορισμό του ορίου\pause. Ιδού:
      \begin{block}{Ορισμός ορίου}
            Έστω μια συνάρτηση ορισμένη σε διάστημα της μορφής $(α,x_0)\cup (x_0,β)$. Λέμε ότι η συνάρτηση τείνει στο $λ\in\mathbb{R}$ καθώς το $x$ τείνει στο $x_0$ όταν: \pause
            \newline
            \newline
            Για κάθε $\epsilon>0$ υπάρχει $δ>0$ ώστε, για κάθε $x\in (α,x_0)\cup (x_0,β)$ με $0<|x-x_0|<δ$ να ισχύει $|f(x)-λ|<\epsilon$
      \end{block}

      \begin{tikzpicture}[remember picture, overlay]
            \draw node at (current page.center) {\includegraphics[width=.8\textwidth]{images/6r7k74}};
      \end{tikzpicture}
\end{frame}

\section{Ασκήσεις}
\begin{frame}{Εξάσκηση}
      Μόνο από το βιβλίο, μόνο γραφικά!
\end{frame}

\section{}
\begin{frame}
      Στο moodle θα βρείτε τις ασκήσεις που πρέπει να κάνετε, όπως και αυτή τη παρουσίαση
\end{frame}

\end{document}
